\documentclass[10pt,a4paper]{amsart}
\usepackage[margin=19mm]{geometry}
\usepackage{amsmath,amssymb,mathtools}
\usepackage[scr=rsfs,cal=euler]{mathalfa}
\usepackage{xfrac}
\usepackage[unicode,hidelinks,bookmarks=false]{hyperref}
\newcommand{\ie}{i.e.\ }
\newcommand{\cf}{cf.\ }
\newcommand{\resp}{respectively\ }
\newcommand{\Subsection}{\subsection}
\providecommand{\nobreakdash}{\nobreak\hbox{-}}
\setlength{\emergencystretch}{2em}
\multlinegap0pt
\usepackage{enumitem}
\usepackage{cleveref}
\def\acts{\curvearrowright}
\newcommand{\nn}[0]{S}
\newcommand{\C}[0]{\mathcal{C}}
\newcommand{\D}[0]{\mathcal{D}}
\newcommand{\fps}[1]{[#1]^{<\omega}}
\newcommand{\orq}[0]{\equiv}
\renewcommand\leq{\leqslant}
\renewcommand\geq{\geqslant}
\theoremstyle{plain}
\newtheorem{theorem}{Theorem}[section]
\newtheorem{corollary}[theorem]{Corollary}
\newtheorem{lemma}[theorem]{Lemma}
\newtheorem{prop}[theorem]{Proposition}
\newtheorem{fact}[theorem]{Fact}
\theoremstyle{remark}
\newtheorem{remark}[theorem]{Remark}
\newtheorem{notation}[theorem]{Notation}
\newtheorem{observation}[theorem]{Observation}
\title{Minimal and intrinsic topologies on monoids of elementary embeddings}
\author{J. de la Nuez, Gonz\u00e1lez, Zaniar Ghadernezhad, Paolo Marimon and Michael Pinsker}
\date{Source: arXiv:2603.28419v3 (CC BY 4.0); the authors' own native \LaTeX{} source, set here in the editorial AMS \texttt{amsart} wrapper (the source's own class is \texttt{article}, which is not redistributed).}
\begin{document}
\maketitle
\noindent\textit{Editorial scope.} The counted claim of this unit is the article's Lemma 5.7 (source label \texttt{skipping terms}), the shortening of a chain of algebraically closed sets with controlled normality condition, delivered together with the authors' complete original proof. Retained uncounted support, in source order, is the whole source-local core that the counted claim needs: the passage of Section~5 that precedes it, including the definitions of chains of algebraically closed sets, of the normality monoid of a chain and of the two intersection conditions that the lemma uses, that is source0.tex lines 975--1089. Every retained passage is a verbatim copy of the authors' own source lines and every cross-reference inside them resolves inside the delivered document. Printed numbering: the authors' own theorem-environment declarations are carried unchanged, with the lemma sharing the counter of the theorem environment and being numbered by section as in the authors' preamble, and the counters are advanced so that the delivered PDF prints the counted claim as Lemma~5.7, exactly as the published article does. Omitted: the title block, abstract and introduction of the article, the parts of the text that precede line 975 and everything after the counted proof. Nothing of the counted statement or of the counted proof is omitted. Class substitution: the source class is the standard \texttt{article} class with the authors' own packages, including \texttt{cleveref} for the cross-references used inside the retained passage and \texttt{enumitem} for the enumerated cases of the statement; the wrapper is AMS \texttt{amsart} carrying the authors' own macro and theorem-environment declarations. Reference handling: no retained passage contains a citation; the article's bibliography is a separate file that is kept verbatim in the eprint directory and is not redistributed, and no reference entry is transcribed, delivered or invented. No mathematics is written, repaired or re-derived here.

\setcounter{section}{5}
\setcounter{theorem}{0}
 \begin{lemma}
 	\label{o:commutation}  For any 
   $C\in\fps{\Omega}$ 
and any $g\in G$, we have $\nn_{g(C)}g= g\nn_{C}$. 
 \end{lemma}
 \begin{proof}
 For the inclusion $\nn_{g(\C)}\,g \subseteq g\,\nn_{\C}$, let
$h \in \nn_{g(\C)}$. Then $s := g^{-1}hg$ stabilises $\C$, that is,
$s \in \nn_{\C}$, and $hg = gs \in g\,\nn_{\C}$. The reverse inclusion is
analogous: given $s \in \nn_{\C}$, set $h := gsg^{-1} \in \nn_{g(\C)}$, so
that $gs = hg \in \nn_{g(\C)}\,g$. Hence $g\,\nn_{\C} \subseteq
\nn_{g(\C)}\,g$, and the two inclusions give
$\nn_{g(\C)}\,g = g\,\nn_{\C}$.
 \end{proof}
 
  \begin{observation}
	   \label{o:chain} Assume that we are given a chain $\C=(\overline{c}_{i})_{i=0}^{k}$, where each $\overline{c}_{i}\in\fps{\Omega}$. Then  $\nn_{\mathcal{C}}$ is precisely the set of $s\in S$ such that there exists a chain %
       ${\mathcal{C}'}:=(\overline{c}_{i}')_{i=0}^{k}$ with $\overline{c}_0'=\overline{c}_0$, $\overline{c}'_{k}=s(\overline{c}_{k})$, and $\overline{c}'_{i}\overline{c}'_{i+1}\orq \overline{c}_{i}\overline{c}_{i+1}$ for every $i<k$.  
  \end{observation}
  \begin{proof}
  Recall that $\nn_{\C} = \nn_{\overline{c}_0}\nn_{\overline{c}_1}\cdots\nn_{\overline{c}_k}$. First assume that $s\in \nn_{\C}$, so $s=g_0\cdots g_k$ with $g_i\in \nn_{\overline{c}_i}$ for each %
  {$i\leq k$.}
  Then set $\overline{c}_i':=g_0\cdots g_i(\overline{c}_i)$ for every $i$. 
  For $i < k$, the element $g_0 \cdots g_i$ witnesses
  $\overline{c}_i\,\overline{c}_{i+1} \orq \overline{c}'_i\,\overline{c}'_{i+1}$: it sends
  $\overline{c}_i$ to $\overline{c}'_i$ by definition, while, using $g_{i+1} \in \nn_{\overline{c}_{i+1}}$, \begin{equation*}
    g_0 \cdots g_i(\overline{c}_{i+1})
      = g_0 \cdots g_i g_{i+1}(\overline{c}_{i+1})
      = \overline{c}'_{i+1}.
  \end{equation*}
Conversely, suppose such a chain 
${\mathcal{C}'}$ exists. For each $i < k$ pick,
  using $\overline{c}_i\,\overline{c}_{i+1} \orq \overline{c}'_i\,\overline{c}'_{i+1}$, an element
  {$g_i \in G$}
  with $g_i(\overline{c}_i\,\overline{c}_{i+1}) = \overline{c}'_i\,\overline{c}'_{i+1}$.
  Set $g_k := s$, so that $g_k(\overline{c}_k) = s(\overline{c}_k) = \overline{c}'_k$ as well. Note that $    s = g_0 \,(g_0^{-1}g_1)\,(g_1^{-1}g_2)\cdots(g_{k-1}^{-1}g_k),
$ and  $g_0 \in \nn_{\overline{c}_0}$ and  $g_{i-1}^{-1}g_i\in \nn_{\overline{c}_i}$ for every $1\leq i\leq k$.  Hence
  $s \in \nn_{\overline{c}_0}\nn_{\overline{c}_1}\cdots\nn_{\overline{c}_k} = \nn_{\C}$.
  

  \end{proof}
  
\begin{observation}\label{obs: subchains}
 Let $\C=(\overline{c}_i)_{i=0}^k$ be a chain, let $J\subseteq \{0, \dots, k\}$ and consider the subchain $\D:=(\overline{c}_i)_{i\in J}$. Then, $\nn_{\D}\subseteq\nn_{\C}$.

\end{observation}
\begin{proof} %
Let $s\in\nn_{\D}$. By definition, $s=\prod_{j\in J} g_j$, {where the $j$'s are in increasing order and} %
$g_j\in\nn_{\overline{c}_j}$. For $i\leq k$, let
\[g'_i:=\begin{cases} g_i &\text{ if } i\in J;\\
1 &\text{ otherwise.}
\end{cases}\]
Then, $s=\prod_{j\in J} g_j=g'_0\cdot \dots \cdot g_k'\in \nn_{\C}$,
concluding the proof. 
\end{proof}


  \begin{corollary}
  	\label{change chains} Assume that $A\in\fps{\Omega}$ and $k\geq 0$. {Let $\C=(\overline{c}_{i})_{i=0}^{k}$  and $\C'=(\overline{c}_{i}')_{i=0}^{k}$ be two chains such that $\overline{c}_{i}\overline{c}_{i+1}\orq_{A} \overline{c}'_{i}\overline{c}'_{i+1}$ for all $0\leq i\leq k-1$.}
    \begin{itemize}
     \item {If $\overline{c}_{0}=\overline{c}'_{0}$, there is $g\in G_{A}$ such that $\nn_{\C'}g=\nn_{\C}$, and $g(\overline{c}'_{k})=\overline{c}_k$;}
        \item {If $\overline{c}_{k}=\overline{c}'_{k}$, there is $g\in G_{A}$ such that $g\nn_{\C}=\nn_{\C'}$ and $g(\overline{c}_{0})=\overline{c}_0'$.}
    \end{itemize}
    
  
  \end{corollary}
  \begin{proof}
  {We argue by induction on the length $k$ of the chains. We begin by proving the first statement. The base case of $k=0$ is trivial since then $\C=\C'$ and we may take $g=\mathrm{Id}$. So, for the inductive hypothesis, suppose the claim holds up to $k\geq 0$ and consider chains $\C=(\overline{c}_{i})_{i=0}^{k+1}$ and $\C'=(\overline{c}'_{i})_{i=0}^{k+1}$ with $\overline{c}_{0}=\overline{c}_{0}'$ and $\overline{c}_{i}\overline{c}_{i+1}\orq_{A} \overline{c}_{i}'\overline{c}'_{i+1}$ for all $0\leq i\leq k$. Denote by $\tilde{\C}$ and $\tilde{\C'}$ the chains obtained from $\C$ and $\C'$ respectively by removing the final term of the chain.}
  
  {By inductive hypothesis, there is $h\in G_A$ such that $\nn_{\tilde{\C'}} h=\nn_{\tilde{\C}}$ and $h(\overline{c}_k')=\overline{c}_k$. Using \cref{o:commutation} for the last equation below, we obtain:}
\begin{equation}\label{eq:chainarg1}
\nn_{\C}=\nn_{\tilde{\C}}\nn_{\overline{c}_{k+1}}=\nn_{\tilde{\C'}}h\nn_{\overline{c}_{k+1}}=\nn_{\tilde{\C'}}\nn_{h^{-1}(\overline{c}_{k+1})}h\;.  
\end{equation}
Since $h\in G_A$,
\[\overline{c}'_k h^{-1}(\overline{c}_{k+1})\equiv_A h(\overline{c}'_k) \overline{c}_{k+1}=\overline{c}_k \overline{c}_{k+1}\equiv_A \overline{c}'_k \overline{c}'_{k+1}\;.\]
{Thus, there is $\ell\in G_{A \overline{c}_k'}$ such that $\ell h^{-1} (\overline{c}_{k+1})=\overline{c}'_{k+1}$. Note that by \cref{o:commutation},} 
\[\nn_{h^{-1}(\overline{c}_{k+1})}=\ell^{-1}\nn_{\ell h^{-1}(\overline{c}_{k+1})} \ell\;.\]
{Furthermore, since $\ell\in G_{\overline{c}'_k}$, $\nn_{\overline{c}_k'} \ell^{-1}=\nn_{\overline{c}_k'}$. (This is just because for $s\in \nn_{\overline{c}_k'}$, we can write $s=s\ell\ell^{-1}$, where $s\ell\in\nn_{\overline{c}_k'}$ too.) Hence, using the two equations above in the right-hand side of \eqref{eq:chainarg1}, we get:}
\[\nn_{\C}=\nn_{\tilde{\C'}}\nn_{h^{-1}(\overline{c}_{k+1})}h=\nn_{\tilde{\C'}}\ell^{-1}\nn_{\ell h^{-1}(\overline{c}_{k+1})} \ell h=\nn_{\tilde{\C'}}\nn_{\overline{c}_{k+1}'} \ell h\;.\]
  {Taking $g:=\ell h$, we can thus conclude the argument for the first part.}

  {The argument for the second part is identical, by carrying out the induction in reverse. That is, now for the inductive step we will denote by $\tilde{\C}$ and $\tilde{\C'}$ the chains obtained from $\C$ and $\C'$ respectively by removing their first term. Then, taking $h\in G_A$ with $h \nn_{\tilde{C}}=\nn_{\tilde{C}'}$ and $h(\overline{c}_1)=\overline{c}_1'$, just as in \eqref{eq:chainarg1}, we obtain} 
  \begin{equation}
    \label{nn equality}	\nn_{\C'}=\nn_{\overline{c}_0'}\nn_{\widetilde{\C'}}=\nn_{\overline{c}_0'}h\nn_{\widetilde{\C}}=h\nn_{h^{-1}(\overline{c}_0')}\nn_{\widetilde{\C}}\;.  
    \end{equation}
{By an analogous argument to the above, we can choose $h'\in G_{A\overline{c}_1}$ such that $h h' (\overline{c}_0)=\overline{c}_0'$, and similarly deduce that} 
\begin{equation*}
    \nn_{\C'}
      = h\,h'\,\nn_{\overline{c}_0}\,\nn_{\widetilde{\C}}
      = h\,h'\,\nn_{\C}.
  \end{equation*}
   {Finally, setting $g := hh'$ yields the desired element of $G_A$.}
  
  \end{proof}
 

  
  
  We remind the reader of the following standard consequence of Neumann's Lemma, which is referred to frequently throughout the paper:  
  \begin{fact}
  	\label{f:neumann} {Let $G\acts\Omega$ with locally finite algebraic closure.} Let $C\in\fps{\Omega}$ be algebraically closed. Then, for all $A, B\in\fps{\Omega}$ there is $A'\in \fps{\Omega} $ such that $A'\equiv_C A$ and $A'\cap B= A\cap B\cap C$.  
  \end{fact}
  
  An easy but useful consequence of this is the following: 
  \begin{lemma}
 	\label{l:intersections} Let $G\acts\Omega$ with locally finite algebraic closure.
 	Then for any algebraically closed $A,C,D\in\fps{\Omega}$  there is some algebraically closed $E\in\fps{\Omega}$ and $g\in G$ such that $E\cap A=C\cap D\cap A$ and the following inclusion holds:
  \begin{equation*}
  	\nn_{E}g\subseteq \nn_{C}\nn_{D}.
  \end{equation*}
 \end{lemma}
 \begin{proof}
 	By \cref{f:neumann}, we can find some element $g\in \nn_{C}\cap G$ such that $g(D)\cap A=C\cap D\cap A$. By applying \cref{o:commutation}, we have $\nn_{g(D)}g=g\nn_{D}\subseteq\nn_{C}\nn_{D}$, so the desired conclusion holds with $E=g(D)$. 
 \end{proof}
 

  \begin{lemma}\label{skipping terms}

        Let $\C=(\overline{c}_i)_{i=0}^{2k}$ be a  chain of algebraically closed sets of length $2k$ for some $k>0$. Assume that we are in one of the following two cases:
	    \begin{enumerate}[label=(\Roman*)]
        \item \label{case1} {There is some algebraically closed $A\in\fps{\Omega}$ such that  $\overline{c}_{2i}\cap \overline{c}_{2i+1}\cap \overline{c}_{2i+2}=A$ for all $0\leq i\leq k-1$.}   
	    	\item \label{case2} 
            {There is some $G$-invariant algebraically closed $\Delta\subseteq\Omega$ such that $\overline{c}_{2i}\cap \overline{c}_{2i+1}\subseteq\Delta$ for all $0\leq i\leq k-1$.} 
            
	    \end{enumerate}
	    Then there is a chain of algebraically closed sets $\D=(\overline{d}_{i})_{i=0}^{k}$ of length $k$ such that for some $g\in G$, 
	    \begin{equation*}
	    	\nn_{\mathcal{D}}{g}\subseteq \nn_{\C}.
	    \end{equation*}

         Moreover: 
	     \begin{itemize}
	    	\item In case \ref{case1} we can ensure $\overline{d}_{i}\cap\overline{d}_{i+1}=A$ for all $0\leq i\leq k-1$, and $g\in G_{A}$;
            \item In case \ref{case2}  we can ensure $\overline{d}_{i}\cap \overline{d}_{i+1}\subseteq\Delta$ for all $0\leq i\leq k-1$.
	    \end{itemize}
  \end{lemma} 

\begin{proof}
{First consider case~\ref{case1}. 
Our aim is to construct a chain of algebraically closed sets
$\C'=(\overline{c}'_{i})_{i=0}^{2k}$ such that 
$\overline{c}'_{0}=\overline{c}_{0}$, 
for all $0\leq i\leq 2k-1$, $\overline{c}_{i}\overline{c}_{i+1}
\equiv_{A}
\overline{c}'_{i}\overline{c}'_{i+1}$, and for all $0\leq i\leq k-1$,
$\overline{c}'_{2i}\cap\overline{c}'_{2i+2}=A$.
Given such a chain of algebraically closed sets $\C'$, we know from the first
part of \cref{change chains} that there is some $g\in G_{A}$ such that $\nn_{\C'}g=
\nn_{\C}$. In particular, letting $\D=(\overline{c}'_{2i})_{i=0}^{k}$, by
Observation~\ref{obs: subchains} we get $\nn_\D\subseteq\nn_{\C'}$, and so $\nn_\D g\subseteq \nn_{\C}$, thus completing the proof.}

{We now construct $\mathcal{C}'$ by induction. The base case is trivial: set $\overline{c}'_{0}=\overline{c}_{0}$. Suppose $(\overline{c}'_{i})_{i=0}^{2l}$ has already been constructed, for some $0\leq l<k$, with all the desired properties; we build $\overline{c}'_{2l+1}$ and $\overline{c}'_{2l+2}$.}

{First choose any $\overline{c}'_{2l+1}$ with $\overline{c}'_{2l}\overline{c}'_{2l+1}
\equiv_{A}
\overline{c}_{2l}\overline{c}_{2l+1}$. We claim there is $\overline{c}'_{2l+2}$ with $\overline{c}'_{2l+1}\overline{c}'_{2l+2}
\equiv_{A}
\overline{c}_{2l+1}\overline{c}_{2l+2}$
and
$\overline{c}'_{2l+2}\cap \overline{c}'_{2l}=A$. To see this, take $\overline{c}''_{2l+2}$ such that $\overline{c}'_{2l}\overline{c}'_{2l+1}
\overline{c}''_{2l+2}
\equiv_A
\overline{c}_{2l}\overline{c}_{2l+1}\overline{c}_{2l+2}$.
From this, and our assumption on $\C$, we get $\overline{c}'_{2l}\cap\overline{c}'_{2l+1}
\cap\overline{c}''_{2l+2}=A$. 
Now apply \cref{f:neumann} to
$\overline{c}''_{2l+2}$,
$\overline{c}'_{2l+1}$, and
$\overline{c}'_{2l}$ to obtain 
$\overline{c}'_{2l+2}\equiv_{A \overline{c}'_{2l+1}}
\overline{c}''_{2l+2}$ with}
\[
\overline{c}'_{2l+2}\cap \overline{c}'_{2l}
=
\overline{c}''_{2l+2}\cap
\overline{c}'_{2l+1}\cap
\overline{c}'_{2l}
=A,
\]
{as required. This completes the construction of $\C'$ in case~\ref{case1}.}

{In case~\ref{case2}, we construct $\C'$ in exactly the same way, but
requiring only that $\overline{c}_{i}\overline{c}_{i+1}
\equiv
\overline{c}'_{i}\overline{c}'_{i+1}$. In the inductive step, having found $\overline{c}_{2l+2}''$ such that $\overline{c}'_{2l}\overline{c}'_{2l+1}
\overline{c}''_{2l+2}
\equiv
\overline{c}_{2l}\overline{c}_{2l+1}\overline{c}_{2l+2}$, we use \cref{f:neumann} to
find $\overline{c}'_{2l+2}\equiv_{\overline{c}'_{2l+1}}
\overline{c}''_{2l+2}$,}
\[
\overline{c}'_{2l}\cap\overline{c}'_{2l+2}=\overline{c}''_{2l+2}\cap
\overline{c}'_{2l+1}\cap
\overline{c}'_{2l}\subseteq\Delta\;,
\]
{which yields the desired containments for $\D$.}
\end{proof}

\end{document}
